\fontsize{11truept}{11truept}\selectfont
\newgeometry{top=45truemm,right=24truemm,left=24truemm}
\thispagestyle{empty}
    % 受験番号タイトル設定日付
    \begin{center}
        % 氏名
        \leavevmode
        \scalebox{1.4}{ ${{\setBold[0.01]令~\;和~\;\exey{~\;}年~\;度~\;京~\;都~\;大~\;学~\;模~\;擬~\;試~\;験}}$ }
        \\[18mm]
        　　　　\scalebox{1.3}{{\huge $数\qquad{学}\quad{(理系)}$}}　　　
        \\[14mm]
        {\large 　$\maxscore$点満点}
        \\[2mm]
        <配点は，一般選抜学生募集要項に記載の通り．って本家は言ってる>
    \end{center}
    % 注意事項
    \leavevmode\\[0mm]
        {\Large 　\;\;$\textgt{{\mdseries (注　意)}}$\;\;　}
    \begin{enumerate}[parsep=0ex, itemsep=0ex, leftmargin=2cm]
        \item[1.\ ]問題冊子および解答冊子は監督者の指示があるまで開かないこと．
        \item[2.\ ]解答冊子は表紙のほかに，解答用ページ，計算用ページ，余白ページをあわせて16ページある．
        \item[3.\ ]問題は全部で$6$題ある($1$ページから$2$ページ) 
        \item[4.\ ]試験開始後，解答冊子の表紙所定欄に学部名・受験番号・氏名をはっきり記入すること．表紙には，それ以外のことを書いてはならない．
        \item[5.\ ]\underbar{解答は問題番号に対応する解答用ページに書くこと．それ以外のページに書かれた}\\\underbar{物は採点の対象としない}．ただし，解答用ページに続き方をはっきり示した場合は，見開きの隣接する計算用ページに「計算用ページに続く」旨が明示されてたときに限って，計算用ページに書かれているものを解答の一部として採点する．なお，他のページに書かれたものは採点の対象とはならないので注意すること． 
        \item[6.\ ]解答に関係ないことを書いた答案は無効にすることがある．なお，計算用紙ページおよび余白ページに書かれた解答のための下書き，解答などは，消さずに残しておいてもよい．
        \item[7.\ ]解答冊子は，どのページも切り離してはならない．
        \item[8.\ ]問題冊子は持ち帰ってもよいが，解答冊子は持ち帰ってはならない．
    \end{enumerate}
